% !TeX root = ../main.tex

\chapter{緒論}

\section{研究背景}

在當前跨領域科學研究中，高通量觀測與影像資訊自動化已成為微觀尺度分析的核心趨勢。其中，倒立式顯微鏡憑藉其靈活的樣本載放空間，成為細胞分析與材料表徵不可或缺的基石。然而，現有市場環境中，高效能螢光成像設備長期受限於高昂的購置成本與高度封閉的黑盒子架構，不僅限制了研究場域的普及性，也阻礙了硬體底層與特殊運算演算法的整合開發。

儘管近年來學界嘗試推動各類開源硬體顯微方案以降低成本，但這類方案往往在機械穩定性上產生妥協。特別是基於純 3D 列印支撐的機構，常因材料特性導致在大行程掃描與長時間定位時出現剛性不足、潛變偏移等物理缺陷。此外，低成本影像感測器與顯微光學系統之間的非理想匹配，亦會引入如亮度不均與色彩偏差等系統性誤差。因此，如何結合高剛性的機械結構、模組化設計邏輯以及運算成像演算法，構建一個兼具成本優勢、自動化控制能力與高解析成像品質的整合型平台，成為當前低成本科研設備研發的重要課題。

\section{研究目的}

本研究旨在研發一套基於模組化結構與智慧化控制的「新型自動化螢光倒立式顯微成像平台」。不同於現有低成本系統，本研究強調在維持經濟性的前提下，顯著提升機械精度與操作自動化。其主要研發目標與技術特徵如下：

首先，在光學基座上導入具備極高幾何公差與標準化特性的樂高（LEGO）組件，構建核心光路架構，並透過高精度 3D 列印件進行光學元件的緊密耦合；其次，針對定位系統，研發一套具備高承載剛性的自動化位移載台，採用 3D 列印機構結合工業級線性滑軌與 T8 導螺桿，解決傳統開源系統在大範圍行程下的抖動與穩定性問題；再者，實作具備動態對焦能力的控制核心與跨平台網頁式圖形介面（Web-based GUI），實現自動化批次掃描與遠端實時觀測功能；最後，整合空間變異校正演算法與圖案化照明傅立葉疊層成像（piFP）技術，從軟體層面修補硬體缺陷並突破物鏡的物理繞射極限。

\section{研究方法}

本研究的研製流程主要分為物理系統開發、演算法管線整合及整體系統確效三個階段。

在物理開發階段，本研究透過標準化模組建立核心光軸，並搭配 NEMA 17 步進馬達與精密機械導件構建剛性龍門位移系統。軟體層面則基於現代開發架構實作了包含自動對焦狀態機、色彩解混張量以及超解析重建在內的完整運算管線。最終，系統驗證階段將利用標準解析度測試板與螢光校正樣本進行量化分析，評估系統在定位精度、自動對焦效率以及多通道螢光成像保真度上的表現，藉以驗證此平台在應對複雜自動化掃描任務時的穩定性與實用價值。
